%document class
\documentclass[10pt,oneside]{book}
%%%% Page Info + Commands %%%%%{

%packages
\usepackage{pgfplots}
\pgfplotsset{compat=1.18}
\usepackage{amsmath}
\usepackage{amsfonts}
\usepackage{amsthm,letltxmacro}
\usepackage{amssymb}
\usepackage{enumerate}
\usepackage{enumitem}
\usepackage[dvipsnames]{xcolor}
\usepackage{changepage}

%This will shrink the page
\newcommand{\smallpage}{  \setlength \oddsidemargin{.5in}
            \setlength \textwidth{5in}
            \setlength \topmargin{0in}
            \setlength \textheight{9in}}


% Commands for sets of numbers
\newcommand{\Z}{\mathbb{Z}}\newcommand{\R}{\mathbb{R}}\newcommand{\N}{\mathbb{N}}

% gordie spacing and boxing commands
\newcommand{\gspc}{\vspace{0.13cm}} % 0.5 space
\newcommand{\gline}{\gspc\\} % 1.5 space
\newcommand{\gbline}{\gspc\gspc\\} % double space
\newcommand{\gbox}[1]{\fbox{\parbox{\linewidth}{#1}}} % multiline box command


\newcommand{\gimage}[3]{
\begin{figure}[h]   
    \centering          
    \includegraphics[width=#2\textwidth]{#1}  
    \caption{#3}
    \label{fig:#1}
\end{figure}\\
}

\newcommand{\centerit}[1]{{\begin{center} #1 \end{center}}}
\newcommand{\ul}[1]{\underline{#1}}

%%%%%%%% command for graphics %%%%%%%%%%%%%
\usepackage{fancyhdr}
%}

%%%% Page Setup %%%%%%%
\smallpage\sloppy
\pagestyle{fancy}
\lhead{\large{\textbf{Day 3 - Berliner lectures us hooray}}}
%\rfoot{\thepage/\pageref{LastPage} }
\setlength{\headheight}{10pt}


% Proof writing
\LetLtxMacro\oldproof\proof
\let\endoldproof\endproof
\renewenvironment{proof}[1][\proofname]
  {\vspace{0.2cm}\begin{adjustwidth}{2em}{2em}
   \oldproof[#1]}
  {\endoldproof
   \end{adjustwidth}}



\begin{document}

\newcommand{\cg}[1]{\color{OliveGreen}#1\color{white}}
\newcommand{\cb}[1]{\color{BlueViolet}#1\color{white}}
\subsection*{Section 2.4}
\begin{enumerate}
	\item [2.] Use the Euclidean Algorithm to obtain integers $x$ and $y$ that satify:
	\begin{enumerate}
		\item [(b)] gcd(24, 138$) = 24x + 138y$.\gline
		We will use the Euclidean algorithm to solve this problem.
		\begin{align*}
			138 &= 5(24) + 18\\
			24 &= 1(18) + 6\\
			18 &= 3(6) + 0
		\end{align*}
		Thus, our gcd is 6. Now we work backwards to find the linear combination.
		\begin{align*}
			6 &= 24 - 1(18)\\
			6&= 24 - 1(138 - 5(24))\\
			6&= 6(24) - 138
		\end{align*}
		Thus, our linear combination will have $x = 6$ and $y = -1$. 
		\item [(d)] gcd(1769, 2378$) = 1769x + 2378y$.\gline
		We will again use the Euclidean algorithm to solve this problem:
		\begin{align*}
			2378 &= 1(1769) + 609\\
			1769 &= 2(609) + 551\\
			609 &= 1(551) + 58\\
			551 &= 9(58) + 29 \\
			58 &= 2(29) + 0
		\end{align*}
		Thus, our gcd is 29. Now, again, we work backwards to find the linear combination:
		\begin{align*}
			29&= 551 -9(58)\\
			29&= 551 - 9(609 - 551)\\
			29 &= 10(551) - 9(609)\\
			29&= 10(1769 - 2(609)) - 9(609)\\
			29 &= -29(609) + 10(1769)\\
			29 &= -29(2378 - 1769) + 10(1769)\\
			29 &= -29(2378) + 39(1769)
		\end{align*}
		This gives us a final result of $x = 39$ and $y = -29$.
	\end{enumerate}
	\item [9.] Prove that the greatest common divisor of two positive integers divides their least common multiple.\gline
	We will show that because the greatest common divisor divides both numbers ($a$ and $b$) in a gcd, it must divide any of their multiples as well. 
	\begin{proof}
		Let $a, b$ be integers with $d = \text{gcd}(a,b)$ and $\ell = \text{lcm}(a,b)$. We want to show that $d\mid \ell$. \gline
		Because $d$ is the greatest common divisor of $a$ and $b$, let $x, y \in \Z$ such that that:
		\begin{align*}
			xd &= a, \text{ and}\\
			yd &= b.
		\end{align*}
		Furthermore, because $\ell$ is the least common multiple of $a$ and $b$, let $ p,q \in \Z$ such that:
		\begin{align*}
			pa &= \ell, \text{ and} \\
			qb &= \ell.
		\end{align*}
		Thus, via substitution, we have that:
		\begin{align*}
			xp\cdot d &= \ell.
		\end{align*}
		Let $m \equiv xp \in \Z$. Hence, because $md = \ell$, via the definition of a divisor, $d=\text{gcd}(a,b)$ divides $\ell = \text{lcm}(a, b)$, as desired. 
	\end{proof}
\end{enumerate}

\subsection*{Section 2.5}
\begin{enumerate}
	\item Which of the following \textbf{Diophantine equations} cannot be solved?\gline
	To solve each of these, we will just find the gcd and determine if the other side of the equation is a multiple of the gcd. 
	\begin{enumerate}
		\item $6x +51y = 22$\gline
		Consider that the gcd of 6 and 51 is equal to 3:
		\begin{align*}
			51 &= 8(6) + \fbox{3}\\
			6 &= 2(3) + 0
		\end{align*}
		Thus, because $22/3$ is not an integer, this Diophantine equation cannot be solved.
		\item $33x + 14y = 115$\gline
		Consider that the gcd of 33 and 14 is 1:
		\begin{align*}
			33 &= 2(14) + 5 \\
			14 &= 2(5) + 4\\
			5 &= 1(4) + \fbox{1}\\
			4 &= 4(1) + 0
		\end{align*}
		Thus, because 115 is trivially a multiple of 1, this Diophjantine equatino can be solved.
		\item $14x + 35y = 93$\gline
		Consider that the gcd of 35 and 14 is 7:
		\begin{align*}
			35 &= 3(14) + \fbox{7}\\
			14 &= 2(7) + 0 
		\end{align*}
		Thus, because $93/7$ is not an integer (as in 7 is not a divisor of 93), this Diophantine equation is not solvable.
	\end{enumerate}
	\item Determine all solutions in the integers of the following Diophantine equations.\gline
	To solve this problem, we will first find the gcd of both $a$ and $b$, and then use our handy-dandy equations:
	\begin{align*}
		x_1 &= x_0 + \left(\frac{b}d\right)t\\
		y_1&= y_0 - \left(\frac{a}d\right)t
	\end{align*}
	To find our solution sets!
	\pagebreak
	\begin{enumerate}
		\item $56x + 72y=40$\gline
		First, let's find the gcd of both numbers:
		\begin{align*}
			72 &= 1(56) + 16\\
			56 &= 3(16) + \fbox 8\\
			16 &= 2(8) + 0
		\end{align*}
		Now that we've identified our gcd as 8 (and that $8\mid 40$), we start looking for our $x_0$ and $y_0$ by doing the Euclidean backwards:
		\begin{align*}
			8&= 56 - 3(16)\\
			8&= 56 - 3(72 - 56)\\
			8&= 4(56) - 3(72)
		\end{align*}
		Thus, we have that $x_0 = 4$ and $y_0 = -3$. 
		\begin{align*}
			x_1 &= 4 + \left(\frac{72}8\right)t &x_1 = 4+9t\\
			y_1&= -3 + \left(\frac{56}8\right)t &y_1 = -3 - 7t
		\end{align*}
		However, these are not our final equations, as they are the solution for $c = 8$. Now, we need to multiply the non-coefficients by $40/8 = 5$ to get our final equations:
		\begin{align*}
			x_1 &= 20+9t\\
			y_1 &= -15 - 7t
		\end{align*}
		\item $24x + 138y = 18$\gline
		First we need to find the gcd the $a$ and $b$:
		\begin{align*}
			138 &= 5(24) + 18\\
			24 &= 1(18) + \fbox 6\\
			18 &= 6 + 0
		\end{align*}
		Now that we have that our gcd is 6, and $6$ clearly divides 18 (so we have solutions), we work backwards to find our particular solution for $x_0$ and $y_0$:
		\begin{align*}
			6&= 24 - 1(18)\\
			6&= 24 - (138 - 5(24))\\
			6&= 6(24) - 138
		\end{align*}
		Thus, we find that $x_0 = 6$ and $y_0 = -1$. Now, we can plug these values into our equations to get a tasty result:
		\begin{align*}
			x_1 &= 6 + \left(\frac{138}6\right)t &x_1 = 6+23t\\
			y_1&= -1 - \left(\frac{24}6\right)t &y_1 = -1 - 4t
		\end{align*}
		Multiplying by our factor of $18/6 = 3$ gives:
		\begin{align*}
			x_1 &= 18+23t\\
			y_1 &= -3 - 4t
		\end{align*}
		Hooray!
		\item Determine all solutions in the positive integers of the following Diophantine equations.\gline
		Little preface from me, gordie! I'm not sure how to solve this one, so I'm going to solve it like I solved the rest of them except restrict $t$ such that by $x_1$ and $y_1$ are positive. 
		\begin{enumerate}
			\item [(b)] $54x + 21y = 906$\gline
			First, let's find the greatest common divisor of both $a$ and $b$:
			\begin{align*}
				54 &= 2(21) + 12\\
				21 &= 1(12) + 9 \\
				12 &= 1(9) + \fbox 3 \\
				9 &= 3(3) + 0
			\end{align*}
			Our gcd is 3, and 3 divides 906 so we have integer solutions to this equation! Solving for $x_0$ and $y_0$ yields:
			\begin{align*}
				3&= 12 - 9\\
				3&= 12 - (21 - 12)\\
				3&= 2(12) - 21\\
				3&= 2(54 - 2(21)) -21 \\
				3&= 2(54) - 5(21)
			\end{align*}
			This gievs us $x_0 = 2$ and $y_0 = -5$. Plugging into our fancy equations gives us that:
			\begin{align*}
				x_1 &= 2 + 7t\\
				y_1 &= -5 - 18t
			\end{align*}
			Multiplying by (906/3) gives:
			\begin{align*}
				x_1 &= 2(302) + 7t\\
				y_1 &= -5(302) - 18t
			\end{align*}
			Now we start to restrict $t$ with $x_1 > 0$ and $y_1>0$:
			\begin{align*}
				t &> -\frac{2(302)}{7}\\
				t &< -\frac{5(302)}{18}
			\end{align*}
			Simplifying to integer forms, $t \ge -86$ and $t <- 83$.
			Thus, our three solutions occur when $t = \{-86, -85, -84\}$
			Writing them out explicitly, we have that:
			\begin{align*}
				x_n &= \{16, 9, 2\}\\
				y_n &= \{2, 20, 38\}
			\end{align*}
			With each solution paired respectively (e.g. (16, 2))
			\pagebreak
		\end{enumerate}
	\end{enumerate}
	\item[6.] A farmer purchased 100 head of livestock for a total cost of \$4000. Prices were as follow: calves, \$120 each; lambs, \$50 each; piglets, \$25 each. If the farmer obtained at least one animal of each type, how many of each did he buy?\gline
	So honestly I'm a little confused by this question, so I think I'm just going to solve it with a system of equations. Let $x$ be the number of calves, $y$ be the number of lambds, and $z$ be the number piglets.
	\begin{align*}
		100 &= x + y + z\\
		4000&= 120x + 50y + 25z
	\end{align*}
	Next, I'm going to reduce this system to two variables so we can use the methods taught in this chapter of the textbook to solve it.
	\begin{align*}
		z&= 100 - x - y, \text{ and then we substitute:} \\
		4000 &= 120 x + 50y + 25(100 - x - y)\\
		1500 &= 95x +25y.
	\end{align*}
	Now, we have an equation of the form $c = ax + by$, which is exactly what we want. Now we use the Euclidean Algorithm to solve for our gcd:
	\begin{align*}
		95 &= 3(25) + 20\\
		25 &= 1(20) + \fbox 5\\
		20 &= 4(5) + 0
	\end{align*}
	Thus, we have a gcd of 5. Now we can do the reverse Euclidean algorithm thing to get our exact values:
	\begin{align*}
		5&= 25 - 20\\
		5 &= 25 - (95 - 3(25))\\
		5 &= 25(4) - 95(1)
	\end{align*}
	We have that  $x = -1$ and $y = 4$. However, this is obviously not realistic as the farmer must buy a minimum of 1 of each livestock. Thus, we must utilize our fancy equations to determine all the possibilities where $x\ge 1$, $y \ge 4$, and $x+y \le 99$ (because have a total of 100 livestock). \gline
	Doing this yields our equations:
	\begin{align*}
		x_1 &= -1 + \frac{25}5t\\
		y_1 &= 4 - \frac{95}5t
	\end{align*} 
	Now, we multiply our constant non-coefficients in the equation by $1500/5 = 300$ and get:
	\begin{align*}
		x_1 &= -300 + 5t\\
		y_1 &= 1200 - 19t
	\end{align*}
	Thus we see that for $x_1> 0$ and $y_1>0$ we must have:
	\begin{align*}
		t &> 60\\
		t &\le 63
	\end{align*}
	So that leaves us with three possibilities: 61, 62, and 63. But this is not satisfactory as we must remain with a singular answer. Thus, we now must satify that final condition, that $x+y \le 99$.
	\pagebreak
	Plugging in our values for t gives us that for $t = \{61, 62, 63\}$ we have that:
	\begin{align*}
		x &= \{5, 10, 15\}\\
		y &= \{41, 22, 3\}
	\end{align*} 
	Turns out all three solutions are valid under this constraint.\\
	Thus, I have create a \textit{LIVESTOCK TABLE} to congratulate me on completing this problem:
	\begin{center}
		\begin{tabular}{c|c|c|c}
			calf & lamb & pig & total\\\hline
			5 & 41 & 54 & 100\\\hline
			10 & 22 & 68 & 100\\\hline
			15 & 3 & 82 & 100\\\hline
		\end{tabular}\gline
		\ul{\textbf{Chart of Total Animal Counts}}\gbline
		\begin{tabular}{c|c|c|c}
			calf & lamb & pig & total\\\hline
			\$600 & \$2050 & \$1350 & \$4000\\\hline
			\$1200 & \$1100 & \$1700 & \$4000\\\hline
			\$1800 & \$150 & \$2050 & \$4000\\\hline
		\end{tabular}\gline
		\ul{\textbf{Chart of Total Animal Costs}}
	\end{center}
	Thus we have a total of 3 possible solutions to the farmer's problem. Which is kinda weird cause the problem made it seem like there was going to be one. 
\end{enumerate}


\end{document}
